\documentclass[a4paper]{article}
\usepackage[pdftex]{graphicx}
\usepackage[utf8]{inputenc}
\usepackage{enumerate}
\usepackage{siunitx}
\sisetup{locale=DE}
\usepackage{href-ul}
\hypersetup{
	colorlinks=true,
	linkcolor=blue,
	urlcolor=blue}
\usepackage{icomma}
\usepackage{amssymb}
\usepackage{geometry}
\geometry{a4paper, top=15mm, left=15mm, right=15mm, bottom=15mm,
	headsep=10mm, footskip=12mm}

\begin{document}
	%Title
	{\bf Oscillating circuit exercise sheet }\\
	\begin{enumerate}[1.]
		%Task 1
		\item A coil with inductance $L$ and a capacitor with capacitance $C$ form a resonant circuit. The ohmic resistance of the resonant circuit is negligibly small. The total energy of the oscillating circuit is $W_{ges}=\SI{2.9e-3}{\joule}$. For the current strength
		$I$ in the oscillatory circuit at a time $t$ with $ t\ge 0~s$ applies:
		$ I(t)= -\hat{I}\cdot \sin{ \left( {2 \pi \over T} \cdot t \right)}$, where $\hat{I}=\SI{120} {\milli\ampere}$ is the peak value of the current and $T=\SI{2.0}{\milli\second}$ is the period of the electromagnetic oscillation.
		
		\begin{enumerate}[a)]
			\item Show the time course of the current $I$ for $0 \le t \le \SI{2,0}{\milli\second}$ in a $t-I$ diagram.
			\item During electromagnetic oscillation, the oscillating circuit capacitor is constantly being charged and discharged. Specify a time interval during which the capacitor is charged. Justify your answer.
			\item Calculate the inductance $L$ of the coil and the capacitance $C$ of the capacitor. $[\textnormal{Partial result:} \quad L=\SI{0,40}{\henry}]$
			\item In the diagram of subtask 1. a) mark the maximum amount $Q_{max}$ of the charge that the capacitor absorbs. Calculate $Q_{max}$.
			\item Calculate the electrical energy $W_{el,1}$ stored in the capacitor for a time $t_1$ at which the current intensity $I$ in the resonant circuit has the value $I=\SI{90}{\milli\ampere} $ accepts.
		\end{enumerate}
		
		%Exercise 2
		\item An oscillating circuit excites a transmitting dipole $D_S$, whose length $l$ can be changed, to oscillate with the frequency $f=\SI{2.5}{\giga\hertz}$. At the length $l_0$, the transmitting dipole is excited in the fundamental oscillation.
		\begin{enumerate}[a)]
			\item Calculate $l_0$.
			
			The transmitting dipole $D_S$ and an identical receiving dipole $D_E$ are arranged perpendicular to a horizontal plane, with this horizontal plane running through the two dipole centers. The receiving dipole registers the oscillation amplitude $E_0$ and receives the power $P_0$. A Hertzian grid is set up between $D_S$ and $D_E$ in such a way that the grid plane is aligned perpendicular to the straight line through the two dipole centers. The bars are significantly longer than the dipoles $D_S$ and $D_E$. First, the grid bars run parallel to the two dipoles.
			
			\item Describe and justify the effects of the grid on the electromagnetic wave in front of and behind the grid.
			
			\item The grid is now rotated by the angle $\alpha=55^\circ$, where the straight line through the two dipole centers is the axis of rotation. Calculate what percent of the power $P_0$ the dipole $D_E$ is now receiving. Illustrate your considerations with a sketch.
		\end{enumerate}
		
	\end{enumerate}
	
	\fbox{
		\begin{minipage}{0.5\textwidth}
			For the solution, please \href{https://www.okuyakl.de/phys/p13schikL235/le235.pdf}{click here} or scan the QR code.\\
			Additional worksheets are available at
			
			\href{http://www.okuyakl.com}{www.okuyakl.com}
		\end{minipage}
		\hfill
		\begin{minipage}{0.4\textwidth}
			\includegraphics[width=1.5 cm]{../../viecher/zwe03}
			\includegraphics[width=3 cm]{qre235}
			\includegraphics[width=2 cm]{../../viecher/afanticon1}
			
	\end{minipage}}
	
\end{document}